\chapter{Conclusioni: Tradeoff Attaccante / Difensore}
\label{ch:conclusioni}

I capitoli precedenti hanno documentato le tecniche e misurato i loro effetti
sugli strumenti di analisi. Questo capitolo sintetizza: dal punto di vista
dell'attaccante, quale scenario è ottimale in funzione del modello difensivo
del bersaglio; dal punto di vista del difensore, quale strategia copre l'intero
spazio delle tecniche analizzate.

% ── 5.1 Scelta Razionale dell'Attaccante ──────────────────────────────────────

\section{Scelta Razionale dell'Attaccante}

L'efficacia delle tecniche di call stack spoofing dipende dal meccanismo con
cui l'EDR del bersaglio percorre la catena di chiamate. Sia Single Frame
Spoofing che Stack Pivoting sono efficaci quando il walker si basa sulle entry
\texttt{.pdata}: il primo inserisce un frame falso con un gadget \texttt{ntdll}
dotato di entry \texttt{.pdata} valida; il secondo sposta SP su un fake stack
allocato separatamente, interrompendo la catena prima che raggiunga il codice
offensivo.

\paragraph{Walker basato sulla catena x29/x30.}
Se l'EDR percorre la catena dei frame pointer invece di affidarsi a
\texttt{.pdata}, le due tecniche si differenziano. Con Single Frame Spoofing
il registro \texttt{x29} di \texttt{InjectExplorer} punta ancora al frame reale
del chiamante: la risalita arriva a \texttt{stack\_spoof.exe} come top frame,
rendendo l'attribuzione al codice offensivo possibile. Con Stack Pivoting SP
viene sostituito con un fake stack privo di catena \texttt{x29/x30} verso il
codice offensivo: il walker non trova continuità e l'attribuzione non risale
all'eseguibile malevolo.


% ── 5.2 Tradeoff del Difensore ────────────────────────────────────────────────

\section{Tradeoff del Difensore: Rilevamento vs Blocco}

Tutti e tre gli scenari presentano firme rilevabili: catena integra per
Baseline, catena \texttt{x29} che rivela il modulo offensivo per Single-Frame Spoofing,
SP fuori TEB per Stack Pivot. Il limite degli strumenti user-mode non è la
copertura ma la sincronia: al momento del rilevamento l'operazione è già
completata.

\begin{itemize}
  \item \textbf{Strumenti user-mode} — copertura completa; costo basso;
        solo rilevamento post-hoc. Il rischio residuo viene accettato
        consapevolmente.
  \item \textbf{EDR con driver kernel-mode} — blocco sincrono nel contesto
        del thread offensivo; punto di partenza autoritativo
        (\texttt{KTRAP\_FRAME}) non modificabile da user-mode; costo in
        licenze e overhead syscall.
\end{itemize}

Non c'è via di mezzo: o si accetta il rischio residuo o si paga il costo
del blocco sincrono.

\bigskip
\noindent
L'asimmetria è strutturale: l'attaccante sceglie la tecnica in base
all'intelligence sul difensore; il difensore deve coprire l'intera superficie
senza sapere quale tecnica sarà impiegata. Si riduce solo alzando il livello
di privilegio del controllo oltre quello raggiungibile dal codice offensivo.
